% GENERATED FILE. Edit canonical lessons, then run npm run generate:books.
% canonical-source-hash: fnv1a64:ead9d7deee379a7d
% canonical-lessons: SA-C218-sammarjani, SA-C218-kutharah, SA-C218-khanitram, SA-C218-halam, SA-C218-datram

\chapter{Broom, Axe, Spade, Plough, Sickle}
\label{ch:sa-broom-axe-spade-plough-sickle}
\hlchaptermodality{218}

\begin{chapteropening}
\textbf{By the end of this chapter:} \emph{I can say a broom, an axe, a spade, a plough and a sickle.}

Complete the last lesson of chapter 218: I can say a broom, an axe, a spade, a plough and a sickle.
\end{chapteropening}

\section[\texorpdfstring{\sk{सम्मार्जनी} (sammārjanī)}{सम्मार्जनी (sammārjanī)}]{\sk{सम्मार्जनी} (sammārjanī) — a broom}

\label{lesson:SA-C218-sammarjani}



\begin{quote}
Before the new one: say the Sanskrit for equal, then the Sanskrit for uneven.
\end{quote}

\subsection*{You'll want to know: \sk{सम्मार्जनी}}
\textbf{\sk{सम्मार्जनी}} — \emph{sammārjanī} — \textquotedblleft{}a broom\textquotedblright{}.

\begin{grammarlens}[title={What you've built}]
One more word for this chapter.
\end{grammarlens}

\subsection*{Your turn}
\begin{itemize}
\raggedright
  \item \emph{Say it:} \emph{sammārjanī}
  \item \emph{Say it:} \emph{sammārjanī}, once more
  \item \emph{Say it:} say \emph{sammārjanī}
  \item \emph{Recall:} say the Sanskrit for equal, then the Sanskrit for uneven, then say \emph{sammārjanī} again
\end{itemize}

\begin{tcolorbox}[breakable,colback=teal!4,colframe=teal!35!black,title={Before you move on}]
What does \sk{सम्मार्जनी} mean? (\textquotedblleft{}A broom\textquotedblright{}.) Say it once more.
\end{tcolorbox}
\section[\texorpdfstring{\sk{कुठारः} (kuṭhāraḥ)}{कुठारः (kuṭhāraḥ)}]{\sk{कुठारः} (kuṭhāraḥ) — an axe}

\label{lesson:SA-C218-kutharah}



\begin{quote}
Before the new one: say the Sanskrit for uneven, then the Sanskrit for a broom.
\end{quote}

\subsection*{You'll want to know: \sk{कुठारः}}
\textbf{\sk{कुठारः}} — \emph{kuṭhāraḥ} — \textquotedblleft{}an axe\textquotedblright{}.

\begin{grammarlens}[title={What you've built}]
One more word for this chapter.
\end{grammarlens}

\subsection*{Your turn}
\begin{itemize}
\raggedright
  \item \emph{Say it:} \emph{kuṭhāraḥ}
  \item \emph{Say it:} \emph{kuṭhāraḥ}, once more
  \item \emph{Say it:} say \emph{kuṭhāraḥ}
  \item \emph{Recall:} say the Sanskrit for uneven, then the Sanskrit for a broom, then say \emph{kuṭhāraḥ} again
\end{itemize}

\begin{tcolorbox}[breakable,colback=teal!4,colframe=teal!35!black,title={Before you move on}]
What does \sk{कुठारः} mean? (\textquotedblleft{}An axe\textquotedblright{}.) Say it once more.
\end{tcolorbox}
\section[\texorpdfstring{\sk{खनित्रम्} (khanitram)}{खनित्रम् (khanitram)}]{\sk{खनित्रम्} (khanitram) — a spade}

\label{lesson:SA-C218-khanitram}



\begin{quote}
Before the new one: say the Sanskrit for a broom, then the Sanskrit for an axe.
\end{quote}

\subsection*{You'll want to know: \sk{खनित्रम्}}
\textbf{\sk{खनित्रम्}} — \emph{khanitram} — \textquotedblleft{}a spade\textquotedblright{}.

\begin{grammarlens}[title={What you've built}]
One more word for this chapter.
\end{grammarlens}

\subsection*{Your turn}
\begin{itemize}
\raggedright
  \item \emph{Say it:} \emph{khanitram}
  \item \emph{Say it:} \emph{khanitram}, once more
  \item \emph{Say it:} say \emph{khanitram}
  \item \emph{Recall:} say the Sanskrit for a broom, then the Sanskrit for an axe, then say \emph{khanitram} again
\end{itemize}

\begin{tcolorbox}[breakable,colback=teal!4,colframe=teal!35!black,title={Before you move on}]
What does \sk{खनित्रम्} mean? (\textquotedblleft{}A spade\textquotedblright{}.) Say it once more.
\end{tcolorbox}
\section[\texorpdfstring{\sk{हलम्} (halam)}{हलम् (halam)}]{\sk{हलम्} (halam) — a plough}

\label{lesson:SA-C218-halam}



\begin{quote}
Before the new one: say the Sanskrit for an axe, then the Sanskrit for a spade.
\end{quote}

\subsection*{You'll want to know: \sk{हलम्}}
\textbf{\sk{हलम्}} — \emph{halam} — \textquotedblleft{}a plough\textquotedblright{}.

\begin{grammarlens}[title={What you've built}]
One more word for this chapter.
\end{grammarlens}

\subsection*{Your turn}
\begin{itemize}
\raggedright
  \item \emph{Say it:} \emph{halam}
  \item \emph{Say it:} \emph{halam}, once more
  \item \emph{Say it:} say \emph{halam}
  \item \emph{Recall:} say the Sanskrit for an axe, then the Sanskrit for a spade, then say \emph{halam} again
\end{itemize}

\begin{tcolorbox}[breakable,colback=teal!4,colframe=teal!35!black,title={Before you move on}]
What does \sk{हलम्} mean? (\textquotedblleft{}A plough\textquotedblright{}.) Say it once more.
\end{tcolorbox}
\section[\texorpdfstring{\sk{दात्रम्} (dātram)}{दात्रम् (dātram)}]{\sk{दात्रम्} (dātram) — a sickle}

\label{lesson:SA-C218-datram}



\begin{quote}
Before the new one: say the Sanskrit for a spade, then the Sanskrit for a plough.
\end{quote}

\subsection*{You'll want to know: \sk{दात्रम्}}
\textbf{\sk{दात्रम्}} — \emph{dātram} — \textquotedblleft{}a sickle\textquotedblright{}.

\begin{grammarlens}[title={What you've built}]
One more word for this chapter.
\end{grammarlens}

\subsection*{Your turn}
\begin{itemize}
\raggedright
  \item \emph{Say it:} \emph{dātram}
  \item \emph{Say it:} \emph{dātram}, once more
  \item \emph{Say it:} say \emph{dātram}
  \item \emph{Recall:} say the Sanskrit for a spade, then the Sanskrit for a plough, then say \emph{dātram} again
\end{itemize}

\begin{tcolorbox}[breakable,colback=teal!4,colframe=teal!35!black,title={Before you move on}]
What does \sk{दात्रम्} mean? (\textquotedblleft{}A sickle\textquotedblright{}.) Say it once more.
\end{tcolorbox}
